\NSPage{157}%
Corollary~\ref{cor:A.3} makes this matrix uniformly invertible on the compact parameter interval. The
nonlinear terms are exactly quadratic. Lemma~\ref{lem:A.2} therefore solves the five equations with
smooth coefficients whose prescribed finite set of \NSi{NS-F-P157-001}{\eta}-derivatives satisfy the required smallness
bounds. All higher fixed derivatives exist with their own finite bounds. This proves \eqref{eq:B.36},
by first fixing the radius of its small-discrepancy neighborhood from data of the outer profile
and then imposing it on the estimates above. The bumps are supported strictly inside their
patch, so the fields equal the ideal near its right endpoint. All five rows and the common
\NSi{NS-F-P157-002}{\Pi_0} now give identical pressure and functions \NSi{NS-F-P157-003}{Q_s,N_s} by \eqref{eq:B.35}; Lemma~\ref{lem:4.4} propagates this
equality along the outer radial profile.
\end{proof}

\subsection{Order of choices and the completed connection.}\label{sec:B.9}
In this subsection, we collect the
parameter choices and state the completed connection. The matching tolerance in \eqref{eq:B.36}
depends only on the outer data, and the radial transition length \NSi{NS-F-P157-004}{T_{\mathrm{sh}}} in \eqref{eq:B.33} has been
bounded independently of the final amplitude \NSi{NS-F-P157-005}{\mathsf C} by Lemma~\ref{lem:B.7}. These facts allow the
amplitude to be chosen to reduce the moment discrepancies before the last transition widths
are fixed, as summarized in \eqref{eq:B.40}.

\begin{remark}\label{rem:B.9}
Fix any finite list of \NSi{NS-F-P157-006}{\eta}-derivative orders whose bounds are needed to preserve
the profile inequalities and solve the finite moment equations, including the extra parameter
derivative in \eqref{eq:B.35}. The choices above and in Section~\ref{sec:A} can be made in the following order:
\NSFormula{NS-F-P157-007}{display}{B.40}
\begin{equation}
\begin{aligned}
 M_d,\ T_d,\ P_*,\ \lambda,\ h
 &\longrightarrow \text{ tolerance for matching moments},\ j_0
 \longrightarrow \delta_*,\sigma_*,\Lambda\\
 &\longrightarrow (\mathcal B_k),\ T_{\mathrm{sh}}
 \longrightarrow \mathsf C,\ X_R
 \longrightarrow \kappa_0,\ t_1,\ \text{ widths of final transitions}.
\end{aligned}
\tag{B.40}\label{eq:B.40}
\end{equation}
The choices for the outer profile first make
\NSi{NS-F-P157-008}{C_{\mathrm{pre}}\sqrt{\lambda}(1+\log(1/\lambda))} small, and then
\NSi{NS-F-P157-009}{h\ll\lambda},
\NSi{NS-F-P157-010}{h\ll e^{-T_d}}, with \NSi{NS-F-P157-011}{C_{\mathrm{pre}}} already fixed. The tolerance determined by the exterior-profile data
is justified by \eqref{eq:B.37}; the choices of \NSi{NS-F-P157-012}{j_0} and \NSi{NS-F-P157-013}{\Lambda} make the contributions
\NSi{NS-F-P157-014}{j_0} and \NSi{NS-F-P157-015}{C^{\mathrm{nat}}_{k_{\mathrm{nat}}}/\Lambda} to the
axial matching error small enough. Lemma~\ref{lem:B.7} fixes \NSi{NS-F-P157-016}{T_{\mathrm{sh}}} before the final amplitude. The
amplitude then ensures
\NSi{NS-F-P157-017}{p_{1,\mathrm r}+p_{2,\mathrm r}^2/p_{1,\mathrm r}>2+c}, the cone inequalities for the outer profile at
large radius, \NSi{NS-F-P157-018}{X_{\mathrm{sep}}/X_R<e^{-8}}, and the bounds on the inner-interval contributions in \eqref{eq:B.39}.
The final narrow transitions secure the activation comparisons and the remaining errors at
\NSi{NS-F-P157-019}{X=X_i} while preserving the previously established bound. Only after these finite choices
is the radial frequency used for admissible-stress-cone realization selected. Smallness is
required only for the prescribed finite set of \NSi{NS-F-P157-020}{\eta}-derivative bounds. Derivatives of every other
fixed order have finite bounds after the choices are made; radial cutoff derivatives may
contain inverse powers of the final widths.
\end{remark}

\begin{corollary}\label{cor:B.10}
The analytic axis profile of Proposition~\ref{prop:B.2} admits a smooth continuation with
\NSi{NS-F-P157-021}{E>0} for \NSi{NS-F-P157-022}{X>0}, matching the reference outer radial profile at
\NSi{NS-F-P157-023}{\log(X/X_R)=-5}, with exact fields,
pressure, and radial moment functions. It has vanishing leading residual stress through
\NSi{NS-F-P157-024}{X_0}, followed
by an inner collar on which the profile remains analytic in \NSi{NS-F-P157-025}{\eta} and satisfies the admissible stress cone
condition, with the flat factor \eqref{eq:B.30}. It satisfies the strict inequalities of the relaxed cone condition
thereafter up to the matching point. Its finite parameters can be chosen in the order \eqref{eq:B.40}.
\end{corollary}

\begin{proof}
Combine Propositions~\ref{prop:B.5} and~\ref{prop:B.8}, the preservation of
\NSi{NS-F-P157-026}{p_1>2} along the profile \eqref{eq:B.34},
and Corollary~\ref{cor:B.6}.
\end{proof}

\section{Realizing the admissible stress cone}\label{sec:C}

The inner and outer profiles \NSi{NS-F-P157-027}{E,U,\Pi} have been joined with the required pressure datum
and radial moments (Corollary~\ref{cor:B.10}). On part of the joining interval, however, the shear
\NSPage{158}%
\NSi{NS-F-P158-001}{(a,-b_s)} satisfies only the relaxed cone condition \eqref{eq:4.21}. The wave construction requires the
additional inequality \NSi{NS-F-P158-002}{v_s>2} from Lemma~\ref{lem:4.5}. In this appendix, we will enforce it by a rapid
radial variation of the shear, supported away from the axis and the exterior.

The shear depends on radial derivatives \eqref{eq:4.11}, while the cumulative moments
\NSi{NS-F-P158-003}{\mathfrak m=(M,I,J,S,C_p)} depend on profile values \eqref{eq:4.15}. We will use this difference to obtain an
order-one change in shear with a small change in the profiles and their moments: we will
prescribe a periodic family of admissible shears \NSi{NS-F-P158-004}{(a_L,-b_L)} whose mean is the original shear
(Lemma~\ref{lem:C.1}), then integrate it at high radial frequency. A correction on the next reserved
interval listed after \eqref{eq:A.9} will restore all five moments exactly and hence preserve the exterior
fields (Proposition~\ref{prop:C.2}).

We take the profile supplied by Corollary~\ref{cor:B.10}, Proposition~\ref{prop:A.4}, and Proposition~\ref{prop:A.7}.
The matching of the axis profile's radial moments to those of the reference outer profile in
Proposition~\ref{prop:B.8} supplies the regular axis and exact moments required in Lemma~\ref{lem:A.8}. These
results therefore give the fixed input profile \NSi{NS-F-P158-005}{E,U,\Pi} used below. We keep this profile and
all its parameters fixed while carrying out the shear modulation and moment correction in
this appendix. Its leading residual stress \NSi{NS-F-P158-006}{\mathcal T_0} is nonzero precisely on
\NSi{NS-F-P158-007}{(X_a,X_b)}, with endpoints
independent of \NSi{NS-F-P158-008}{\eta}. The profile satisfies the strict inequalities of the relaxed cone condition
throughout that interval and the admissible stress cone condition on an inner collar and
from the intermediate power-law interval onward (Propositions~\ref{prop:A.4}, \ref{prop:A.7} and~\ref{prop:A.10}). The two
edge factorizations are \eqref{eq:B.30} and Equations~\eqref{eq:A.48} to~\eqref{eq:A.49}. We choose a compact interval
\NSi{NS-F-P158-009}{\mathcal I=[X_-,X_+]\Subset(X_a,X_b)} containing every point where the admissible stress cone condition
may fail, starting inside the first collar on which the admissible stress cone condition holds
and ending in the intermediate power-law interval before its first reserved patch. Shorter
collars at \NSi{NS-F-P158-010}{X_a} and \NSi{NS-F-P158-011}{X_b} will remain unchanged. We fix the rectangle near the axis from
Corollary~\ref{cor:B.6}, on which \NSi{NS-F-P158-012}{E/\sqrt{2X},U,V_0/X,\Pi} are smooth in
\NSi{NS-F-P158-013}{X,\eta} and analytic in \NSi{NS-F-P158-014}{\eta} on one
common complex neighborhood. We choose its radial endpoint in the shorter inner collar,
before \NSi{NS-F-P158-015}{X_-}. We denote this endpoint by \NSi{NS-F-P158-016}{X_{\mathrm{an}}}.

\subsection{A loop with the prescribed mean.}\label{sec:C.1}
In this subsection, we will construct the family
\NSi{NS-F-P158-017}{(a_L,-b_L)} of shear vectors at each profile point, holding the integrated residual coefficient
\NSi{NS-F-P158-018}{p_s}
fixed (Lemma~\ref{lem:C.1}). Their average must equal \NSi{NS-F-P158-019}{(a,-b_s)}, so the difference has mean zero and
hence a periodic antiderivative. In Proposition~\ref{prop:C.2} we will integrate this difference into a
small change of the profiles. The family must depend smoothly on \NSi{NS-F-P158-020}{(X,\eta)} and remain equal
to the original shear near the radial boundaries, where no modification is needed.

\begin{lemma}\label{lem:C.1}
Let \NSi{NS-F-P158-021}{a,b_s,p_s} be smooth on \NSi{NS-F-P158-022}{\mathcal I\times[-1,1]}, with
\NSi{NS-F-P158-023}{a>0}, and set
\NSi{NS-F-P158-024}{t_s=-b_s/a}, \NSi{NS-F-P158-025}{v_s=a(1+t_s^2)}.
Suppose the strict inequalities of the relaxed cone condition in \eqref{eq:4.21} hold everywhere, and that the
admissible stress cone holds in neighborhoods of the two radial boundary components. There is a
smooth loop \NSi{NS-F-P158-026}{(a_L,-b_L)(X,\eta,\varphi)}, of period one in \NSi{NS-F-P158-027}{\varphi}, with
\NSi{NS-F-P158-028}{a_L>0}, such that
\NSFormula{NS-F-P158-029}{display}{C.1}
\begin{equation}
 \int_0^1 (a_L,-b_L)\,d\varphi=(a,-b_s).
\tag{C.1}\label{eq:C.1}
\end{equation}
Let \NSi{NS-F-P158-030}{\Psi} be the vector of cone gaps in \eqref{eq:4.35}. There are constants
\NSi{NS-F-P158-031}{\kappa_L,\delta_\partial>0}, depending on the fixed data
\NSi{NS-F-P158-032}{a,b_s,p_s,\mathcal I}, such that
\NSFormula{NS-F-P158-033}{display}{C.2}
\begin{equation}
\begin{gathered}
 \Psi_j\bigl(a_L(X,\eta,\varphi),b_L(X,\eta,\varphi),p_s(X,\eta)\bigr)\geq\kappa_L,\\
 (X,\eta,\varphi)\in\mathcal I\times[-1,1]\times(\R/\Z),
 \qquad j=1,2,3,4.
\end{gathered}
\tag{C.2}\label{eq:C.2}
\end{equation}
\NSPage{159}%
On the radial boundary collars,
\NSFormula{NS-F-P159-001}{display}{C.3}
\begin{equation}
\begin{gathered}
 (a_L,-b_L)(X,\eta,\varphi)=(a,-b_s)(X,\eta),\\
 X\in\mathcal I,\qquad \operatorname{dist}(X,\partial\mathcal I)<\delta_\partial,
 \qquad (\eta,\varphi)\in[-1,1]\times(\R/\Z).
\end{gathered}
\tag{C.3}\label{eq:C.3}
\end{equation}
All smoothness assertions include the closed parameter interval
\NSi{NS-F-P159-002}{-1\leq\eta\leq1}.
\end{lemma}

\begin{proof}
Write a shear vector as a positive multiple of \NSi{NS-F-P159-003}{(1,t)}. We will vary the ratio \NSi{NS-F-P159-004}{t} with
mean \NSi{NS-F-P159-005}{t_s}, choose a scalar \NSi{NS-F-P159-006}{v>2}, and seek admissible vectors of the form
\NSi{NS-F-P159-007}{v(1,t)/(1+t^2)}. The
variance of \NSi{NS-F-P159-008}{t} will determine \NSi{NS-F-P159-009}{v}. A change of parameter around the loop will then give these
vectors the required mean \NSi{NS-F-P159-010}{(a,-b_s)}.

For an auxiliary angle \NSi{NS-F-P159-011}{\theta'}, write
\NSi{NS-F-P159-012}{\langle f\rangle_{\theta'}=(2\pi)^{-1}\int_0^{2\pi}f(\theta')\,d\theta'}; this average is distinct from
the physical angular and auxiliary-torus averages. Put
\NSi{NS-F-P159-013}{P_{\mathrm c}(t)=p_{s,1}+p_{s,2}t} and
\NSi{NS-F-P159-014}{J_{\mathrm c}(t)=p_{s,2}-p_{s,1}t}. Choose
\NSi{NS-F-P159-015}{0<d_0<\tfrac12\min(P_{\mathrm c}(t_s)-2)}, and define
\NSFormula{NS-F-P159-016}{display}{C.4}
\begin{equation}
 M_e(z)=\bigl\langle e^{z\sin\theta'}\bigr\rangle_{\theta'},
\tag{C.4}\label{eq:C.4}
\end{equation}
\NSFormula{NS-F-P159-017}{display}{C.5}
\begin{equation}
 t(\theta';\mu)=t_s+d_0
 \frac{e^{\mu p_{s,2}\sin\theta'}/M_e(\mu p_{s,2})-1}{p_{s,2}},
 \qquad \mu\geq0.
\tag{C.5}\label{eq:C.5}
\end{equation}
The apparent singularity at \NSi{NS-F-P159-018}{p_{s,2}=0} is removable. Indeed, for a smooth numerator
\NSi{NS-F-P159-019}{G(p)} with
\NSi{NS-F-P159-020}{G(0)=0},
\NSi{NS-F-P159-021}{G(p)/p=\int_0^1G'(up)\,du}. Applying this identity to \eqref{eq:C.5} gives joint smoothness in
all its variables and the value \NSi{NS-F-P159-022}{t=t_s+d_0\mu\sin\theta'} when
\NSi{NS-F-P159-023}{p_{s,2}=0}. The family has the identities
\NSFormula{NS-F-P159-024}{display}{none}
\begin{gather*}
 \langle t\rangle_{\theta'}=t_s,
 \qquad P_{\mathrm c}(t)\geq P_{\mathrm c}(t_s)-d_0>2,\\
 V(\mu,p_{s,2}):=\bigl\langle(t-t_s)^2\bigr\rangle_{\theta'}
 =\frac{d_0^2}{p_{s,2}^2}
 \left[\frac{M_e(2\mu p_{s,2})}{M_e(\mu p_{s,2})^2}-1\right].
\end{gather*}
At \NSi{NS-F-P159-025}{p_{s,2}=0} the last expression is \NSi{NS-F-P159-026}{d_0^2\mu^2/2}.

We show that a prescribed finite variance can be attained uniformly over the compact
range of \NSi{NS-F-P159-027}{p_{s,2}}, including \NSi{NS-F-P159-028}{p_{s,2}=0}. Let
\NSi{NS-F-P159-029}{g=\log M_e}. Its second derivative is the variance of
\NSi{NS-F-P159-030}{\sin\theta'}
under the strictly positive tilted probability density, so \NSi{NS-F-P159-031}{g''>0}. For
\NSi{NS-F-P159-032}{p\neq0} and \NSi{NS-F-P159-033}{\mu>0},
\NSFormula{NS-F-P159-034}{display}{C.6}
\begin{equation}
 \partial_\mu\{g(2\mu p)-2g(\mu p)\}
 =2p\{g'(2\mu p)-g'(\mu p)\}>0.
\tag{C.6}\label{eq:C.6}
\end{equation}
Thus \NSi{NS-F-P159-035}{V} increases strictly with \NSi{NS-F-P159-036}{\mu>0}, also when
\NSi{NS-F-P159-037}{p=0}. Quadratic upper and lower bounds
for \NSi{NS-F-P159-038}{\sin\theta'} near its maximum, and a uniform gap from that maximum on its complement, give
\NSi{NS-F-P159-039}{M_e(z)\asymp e^{|z|}/\sqrt{|z|}} for \NSi{NS-F-P159-040}{|z|\geq1}. Consequently
\NSi{NS-F-P159-041}{M_e(2z)/M_e(z)^2\asymp\sqrt{|z|}}, and
\NSi{NS-F-P159-042}{V(\mu,p)\to\infty} for
every fixed \NSi{NS-F-P159-043}{p}, including zero. Continuity, monotonicity, and a finite cover of the compact
range of \NSi{NS-F-P159-044}{p_{s,2}} now give one finite \NSi{NS-F-P159-045}{\mu_{\max}} such that
\NSFormula{NS-F-P159-046}{display}{C.7}
\begin{equation}
 V(\mu_{\max},p_{s,2})>3/\min a
 \qquad\text{throughout }\mathcal I\times[-1,1].
\tag{C.7}\label{eq:C.7}
\end{equation}

We can therefore choose the variance of \NSi{NS-F-P159-047}{t} throughout the compact parameter set. We
next bound the allowed value of \NSi{NS-F-P159-048}{v} so that all resulting shear vectors stay inside the cone.
For \NSi{NS-F-P159-049}{0\leq\mu\leq\mu_{\max}}, the functions
\NSi{NS-F-P159-050}{t(\theta';\mu)} form a compact family with
\NSi{NS-F-P159-051}{P_{\mathrm c}(t)>2}. The upper
admissible bound \NSi{NS-F-P159-052}{\mathcal U} for \NSi{NS-F-P159-053}{v} from \eqref{eq:4.21}, the smaller root in
\NSi{NS-F-P159-054}{v} of
\NSi{NS-F-P159-055}{2(P_{\mathrm c}-v)^2=(v-2)J_{\mathrm c}^2}, is
continuous and strictly greater than two there. Choose \NSi{NS-F-P159-056}{0<\delta_L<1} so small that
\NSFormula{NS-F-P159-057}{display}{C.8}
\begin{equation}
 \mathcal U(P_{\mathrm c}(t),J_{\mathrm c}(t))>2+\delta_L
 \qquad(0\leq\mu\leq\mu_{\max}).
\tag{C.8}\label{eq:C.8}
\end{equation}
Shrink \NSi{NS-F-P159-058}{\delta_L} further below the positive minimum of \NSi{NS-F-P159-059}{v_s-2} on smaller neighborhoods of the
radial boundaries. Choose \NSi{NS-F-P159-060}{\delta_L} after fixing \NSi{NS-F-P159-061}{\mu_{\max}}.

\NSPage{160}%
Only the region where \NSi{NS-F-P160-001}{v_s} is near or below two needs a nonzero variance. We leave the
admissible boundary data unchanged. Let \NSi{NS-F-P160-002}{\zeta_L} be a smooth nonnegative function of
\NSi{NS-F-P160-003}{v_s}, at
most one, equal to one for \NSi{NS-F-P160-004}{v_s\leq2+\delta_L/8} and zero for
\NSi{NS-F-P160-005}{v_s\geq2+\delta_L/4}. Set
\NSFormula{NS-F-P160-006}{display}{C.9}
\begin{equation}
 v_*=2+\delta_L/2,
 \qquad \rho=\zeta_L(v_s)^2(v_*-v_s),
 \qquad v=v_s+\rho,
\tag{C.9}\label{eq:C.9}
\end{equation}
where \NSi{NS-F-P160-007}{\rho} is defined to be zero off the support of the cutoff. Its nonnegative square root
is smooth: wherever the cutoff can be nonzero, \NSi{NS-F-P160-008}{v_*-v_s\geq\delta_L/4}, so it is
\NSi{NS-F-P160-009}{\zeta_L(v_s)\sqrt{v_*-v_s}},
extended by zero. Also \NSi{NS-F-P160-010}{0\leq\rho<3}. Solve
\NSFormula{NS-F-P160-011}{display}{C.10}
\begin{equation}
 V(\mu,p_{s,2})=\rho/a,
 \qquad 0\leq\mu<\mu_{\max}.
\tag{C.10}\label{eq:C.10}
\end{equation}
Existence and uniqueness follow from \eqref{eq:C.7} and strict monotonicity. To verify smooth
dependence where \NSi{NS-F-P160-012}{\rho=0}, the expansion of \eqref{eq:C.4} gives, uniformly for bounded
\NSi{NS-F-P160-013}{p},
\NSFormula{NS-F-P160-014}{display}{none}
\[
 V(\mu,p)=\tfrac12d_0^2\mu^2+O(\mu^4p^2).
\]
The signed square root of \NSi{NS-F-P160-015}{V} is smooth and odd near \NSi{NS-F-P160-016}{\mu=0}, with derivative
\NSi{NS-F-P160-017}{d_0/\sqrt2}. The
implicit function theorem applied to this square root and the smooth right side
\NSi{NS-F-P160-018}{\sqrt\rho/\sqrt a}
proves smoothness through \NSi{NS-F-P160-019}{\rho=0}. Away from zero it follows from \eqref{eq:C.6}.

The definition of \NSi{NS-F-P160-020}{\rho} gives
\NSFormula{NS-F-P160-021}{display}{none}
\[
 v=(1-\zeta_L(v_s)^2)v_s+\zeta_L(v_s)^2v_*.
\]
Since \NSi{NS-F-P160-022}{0\leq\zeta_L\leq1}, the three cutoff regimes satisfy
\NSFormula{NS-F-P160-023}{display}{none}
\[
\begin{aligned}
 v&=v_* &&\text{if }v_s\leq2+\delta_L/8,\\
 2<v_s\leq v&\leq v_* &&\text{if }2+\delta_L/8<v_s<2+\delta_L/4,\\
 v&=v_s>2 &&\text{if }v_s\geq2+\delta_L/4.
\end{aligned}
\]
Thus \NSi{NS-F-P160-024}{v>2} everywhere, and \NSi{NS-F-P160-025}{v\leq v_*} wherever the cutoff is nonzero. Where the cutoff
is zero, \NSi{NS-F-P160-026}{\mu=0} and \NSi{NS-F-P160-027}{v=v_s>2}, so the strict inequalities of the relaxed cone condition
apply to the given data \NSi{NS-F-P160-028}{a,b_s,p_s}. It follows from \eqref{eq:C.8} that all the proposed states satisfy
\NSi{NS-F-P160-029}{2<v<\mathcal U(P_{\mathrm c}(t),J_{\mathrm c}(t))}.

The vectors now lie in the admissible cone. Their average depends on how long the loop
spends at each value of \NSi{NS-F-P160-030}{t}. To impose both prescribed shear averages, reparametrize
\NSi{NS-F-P160-031}{\theta'} by a
lifted angle \NSi{NS-F-P160-032}{\varphi}, with origin at \NSi{NS-F-P160-033}{\theta'=0}, by
\NSFormula{NS-F-P160-034}{display}{none}
\[
 \frac{d\varphi}{d\theta'}=\frac{a(1+t^2)}{2\pi v},
 \qquad (a_L,-b_L)=\frac{v(1,t)}{1+t^2}.
\]
Since \NSi{NS-F-P160-035}{a\langle1+t^2\rangle_{\theta'}=v_s+aV=v}, the lift satisfies
\NSi{NS-F-P160-036}{\varphi(\theta'+2\pi)=\varphi(\theta')+1}. Its derivative has a
positive lower bound on the compact family. It therefore defines a circle diffeomorphism
with smoothly parameter-dependent inverse. Changing variables gives
\NSFormula{NS-F-P160-037}{display}{none}
\[
 \int_0^1a_L\,d\varphi=a,
 \qquad
 \int_0^1(-b_L)\,d\varphi=a\langle t\rangle_{\theta'}=at_s=-b_s.
\]
The loop satisfies \NSi{NS-F-P160-038}{-b_L/a_L=t} and \NSi{NS-F-P160-039}{a_L(1+t^2)=v}, so the admissible stress cone inequalities
follow from the preceding scalar inequalities. On the boundary neighborhoods the cutoff
vanishes, hence the loop equals the original shear \NSi{NS-F-P160-040}{(a,-b_s)}. Compactness gives the uniform
margins \eqref{eq:C.2} in \NSi{NS-F-P160-041}{a_L}, \NSi{NS-F-P160-042}{v-2}, and the two strict inequalities of the quadratic cone test; the fixed
boundary neighborhoods give \eqref{eq:C.3}.
\end{proof}

\NSPage{161}%
\subsection{Radial modulation and exact restoration of the radial moment functions.}\label{sec:C.2}
The loop
in Lemma~\ref{lem:C.1} was constructed with \NSi{NS-F-P161-001}{p_s} fixed. In this subsection, we will insert it into the
profiles \NSi{NS-F-P161-002}{E,U} and control the resulting change in \NSi{NS-F-P161-003}{p_s}, which depends on their radial moments
through \eqref{eq:4.16}. We will evaluate periodic antiderivatives at \NSi{NS-F-P161-004}{N\log X} and divide by \NSi{NS-F-P161-005}{N}, so that
the profile values change by \NSi{NS-F-P161-006}{O(N^{-1})} while the shear changes by order one. The proof of
Proposition~\ref{prop:C.2} will establish these estimates, bound the resulting moment errors by
\NSi{NS-F-P161-007}{O(N^{-1})},
and remove them on the reserved correction interval. The comparison with the fixed input
profile \NSi{NS-F-P161-008}{E,U,\Pi} in the following proposition does not assert small radial derivatives.

\begin{proposition}\label{prop:C.2}
For the fixed input profile \NSi{NS-F-P161-009}{E,U,\Pi} selected at the beginning of this appendix,
a sufficiently large finite integer \NSi{NS-F-P161-010}{N} and a correction supported in the first reserved patch in the
intermediate power-law interval give smooth profiles \NSi{NS-F-P161-011}{\widetilde E,\widetilde U,\widetilde\Pi} with
\NSi{NS-F-P161-012}{\widetilde E>0} for \NSi{NS-F-P161-013}{X>0}, satisfying
the admissible stress cone throughout \NSi{NS-F-P161-014}{(X_a,X_b)}. Both endpoint collars, the rectangle near the axis
described above, and the two patches reserved for higher-order and mean corrections are unchanged.
With \NSi{NS-F-P161-015}{\widetilde{\mathfrak m},\widetilde Q_s,\widetilde N_s} defined from these profiles by \eqref{eq:4.15} and \eqref{eq:4.16},
\NSFormula{NS-F-P161-016}{display}{none}
\[
 (\widetilde{\mathfrak m},\widetilde\Pi,\widetilde Q_s,\widetilde N_s)
 =(\mathfrak m,\Pi,Q_s,N_s)
 \qquad\text{beyond the first correction patch}.
\]
Before that point, profile values and each fixed number of \NSi{NS-F-P161-017}{\eta}-derivatives differ from those of the input
profile by \NSi{NS-F-P161-018}{O(N^{-1})}.
\end{proposition}

\begin{proof}
The prescribed averages make the differences between the loop and the original
shear integrable in the periodic variable. Let \NSi{NS-F-P161-019}{\mathscr A,\mathscr B} be the unique zero-mean periodic
antiderivatives
\NSFormula{NS-F-P161-020}{display}{C.11}
\begin{equation}
 \partial_\varphi\mathscr A=-\tfrac12(a_L-a),
 \qquad
 \partial_\varphi\mathscr B=\tfrac12E(b_L-b_s).
\tag{C.11}\label{eq:C.11}
\end{equation}
They exist by \eqref{eq:C.1}, are smooth, and vanish identically on the boundary neighborhoods of
\NSi{NS-F-P161-021}{\mathcal I}.
Extend them by zero radially and put
\NSFormula{NS-F-P161-022}{display}{C.12}
\begin{equation}
 E_N=E\exp\!\left(\frac{\mathscr A(X,\eta,N\log X)}{N}\right),
 \qquad
 U_N=U+\frac{\mathscr B(X,\eta,N\log X)}{N}.
\tag{C.12}\label{eq:C.12}
\end{equation}
The extension agrees smoothly with the input profiles \NSi{NS-F-P161-023}{E,U}. In the next identity,
\NSi{NS-F-P161-024}{\mathcal D_X=X\partial_X}
differentiates a function of \NSi{NS-F-P161-025}{(X,\eta,\varphi)} while holding \NSi{NS-F-P161-026}{\varphi} fixed; all its terms are then evaluated at
\NSi{NS-F-P161-027}{\varphi=N\log X}. The exact shears of \eqref{eq:C.12} are
\NSFormula{NS-F-P161-028}{display}{C.13}
\begin{equation}
 a_N=a_L-\frac{2\mathcal D_X\mathscr A}{N},
 \qquad
 b_N=e^{-\mathscr A/N}\left(b_L+\frac{2\mathcal D_X\mathscr B}{NE}\right).
\tag{C.13}\label{eq:C.13}
\end{equation}
Indeed \NSi{NS-F-P161-029}{X\partial_X} after substitution is
\NSi{NS-F-P161-030}{\mathcal D_X+N\partial_\varphi}; inserting \eqref{eq:C.11} into
\NSi{NS-F-P161-031}{a_N=1-2X\partial_X\log E_N} and
\NSi{NS-F-P161-032}{b_N=2X\partial_XU_N/E_N} proves both formulas, including the exponential denominator in the axial
shear.

On the fixed compact radial range from \NSi{NS-F-P161-033}{X_-} through the first correction interval,
\NSi{NS-F-P161-034}{X,E,H,L}
have positive lower bounds. For each fixed \NSi{NS-F-P161-035}{m\geq0}, the phase in \eqref{eq:C.12} is independent of
\NSi{NS-F-P161-036}{\eta}, so
\NSFormula{NS-F-P161-037}{display}{C.14}
\begin{equation}
 \sup_X\sum_{j=0}^m
 \left(\left|\partial_\eta^j(E_N-E)\right|
       +\left|\partial_\eta^j(U_N-U)\right|\right)
 \leq C_mN^{-1}.
\tag{C.14}\label{eq:C.14}
\end{equation}
The same estimate holds for the differences between the shears in \eqref{eq:C.13} and their loop
values. The derivative \NSi{NS-F-P161-038}{X\partial_X(E_N-E)} can have order one. In general, for fixed
\NSi{NS-F-P161-039}{r\geq1}, \NSi{NS-F-P161-040}{m\geq0},
radial differentiation of \eqref{eq:C.12} gives a bound \NSi{NS-F-P161-041}{C_{r,m}N^{r-1}} for the mixed derivatives of the profile
difference. To check the cone condition for the modified profiles, we need control of \NSi{NS-F-P161-042}{p_s} as
well as of the shear. Its radial integral formulas use profile values and parameter derivatives,
so the preceding estimates suffice even though radial derivatives need not be small.

\NSPage{162}%
Let \NSi{NS-F-P162-001}{\mathfrak m=(M,I,J,S,C_p)} and \NSi{NS-F-P162-002}{\mathfrak m_N} denote the five radial moment functions in \eqref{eq:4.15} for
\NSi{NS-F-P162-003}{E,U}
and \NSi{NS-F-P162-004}{E_N,U_N}, respectively. Keep the same pressure value \NSi{NS-F-P162-005}{\Pi_0(\eta)} at the axis, so
\NSi{NS-F-P162-006}{\Pi_N=\Pi_0+C_{p,N}}.
Integration of \eqref{eq:C.14} on this fixed radial range, with the bounded weights in the radial
moment definitions, gives
\NSFormula{NS-F-P162-007}{display}{C.15}
\begin{equation}
 \sup_X\sum_{j=0}^m\left|\partial_\eta^j(\mathfrak m_N-\mathfrak m)\right|
 +\sup_X\sum_{j=0}^m\left|\partial_\eta^j(\Pi_N-\Pi)\right|
 \leq C_mN^{-1}.
\tag{C.15}\label{eq:C.15}
\end{equation}
The formulas for \NSi{NS-F-P162-008}{Q_s,N_s} in \eqref{eq:4.16} involve these radial moment functions and their first
parameter derivatives, together with profile values. Their denominators are bounded away
from zero here. Consequently
\NSFormula{NS-F-P162-009}{display}{C.16}
\begin{equation}
 \sup_X\sum_{j=0}^m\left|\partial_\eta^j(p_{s,N}-p_s)\right|
 \leq C_mN^{-1}.
\tag{C.16}\label{eq:C.16}
\end{equation}
The estimate at order \NSi{NS-F-P162-010}{m} uses profile and radial moment estimates through order
\NSi{NS-F-P162-011}{m+1} in \NSi{NS-F-P162-012}{\eta}.
It uses no comparison of radial derivatives. The uniform positive lower bounds in the cone
inequalities of Lemma~\ref{lem:C.1} and \eqref{eq:C.13} now preserve the admissible stress cone on
\NSi{NS-F-P162-013}{\mathcal I} for large
\NSi{NS-F-P162-014}{N}. Between \NSi{NS-F-P162-015}{\mathcal I} and the first correction interval, the fields equal
\NSi{NS-F-P162-016}{E,U} and \NSi{NS-F-P162-017}{p_{s,N}-p_s} satisfies
\eqref{eq:C.16}, so the uniform positive lower bounds in the input profile's admissible stress cone
inequalities persist there as well.

The modified profiles now satisfy the cone inequalities through the joining interval. Their
moments have errors of order \NSi{NS-F-P162-018}{N^{-1}}, which must be removed before reaching the exterior.
We correct them on the first patch listed after \eqref{eq:A.9}; it precedes the terminal compensation
patch and the two later correction patches. On this patch the input profiles satisfy
\NSi{NS-F-P162-019}{U=0},
\NSi{NS-F-P162-020}{E=K(\eta)X^{-1/2-\lambda}}, where \NSi{NS-F-P162-021}{K} is smooth and has a positive lower bound. Add two fixed
compact bumps to \NSi{NS-F-P162-022}{U} and three to \NSi{NS-F-P162-023}{E}, with coefficients depending smoothly on
\NSi{NS-F-P162-024}{\eta}. Order the
five radial moment functions at the right endpoint as \NSi{NS-F-P162-025}{(M,J;I,S,C_p)}. At zero coefficients
their differential is block diagonal, with weights
\NSFormula{NS-F-P162-026}{display}{none}
\[
\begin{array}{c|c|c}
 \text{rows} & \text{weights in }dX & \text{powers of }X\\ \hline
 M,J & 1,H & 0,-\lambda\\
 I,S,C_p & \sqrt{2X},-E,E/X & 1/2,-1/2-\lambda,-3/2-\lambda.
\end{array}
\]
The exponents are distinct for the already fixed \NSi{NS-F-P162-027}{\lambda>0}. Choose the bumps with ordered
disjoint supports within each block. Lemma~\ref{lem:A.1}, or its specialization to these five moment
functions Corollary~\ref{cor:A.3}, gives an invertible matrix, with uniform inverse on
\NSi{NS-F-P162-028}{[-1,1]}. Its
bound can depend on \NSi{NS-F-P162-029}{\lambda} and the fixed patch scales; the two
\NSi{NS-F-P162-030}{U}-weights coalesce when \NSi{NS-F-P162-031}{\lambda\to0},
so no uniformity in that limit is asserted.

The map from correction coefficients to changes in \NSi{NS-F-P162-032}{(M,J,I,S,C_p)} equals this linear map
plus quadratic terms: \NSi{NS-F-P162-033}{J} contains the product of the two field changes, and
\NSi{NS-F-P162-034}{S,C_p} contain
their squares. The discrepancy accumulated before this correction is \NSi{NS-F-P162-035}{O_m(N^{-1})} by \eqref{eq:C.15}.
Lemma~\ref{lem:A.2} therefore gives a smooth solution for the coefficients near zero which cancels that
discrepancy. For every fixed parameter order its coefficients are \NSi{NS-F-P162-036}{O_m(N^{-1})}. Choose
\NSi{NS-F-P162-037}{N} after \NSi{NS-F-P162-038}{\lambda},
the patch scales, and all choices entering the fixed input profile so that the discrepancy lies in
the neighborhood where this solution is defined and its derivatives through the finitely many
\NSi{NS-F-P162-039}{\eta}-orders needed for positivity and the cone inequalities meet their tolerances. Differentiating
the implicit equation gives finite bounds for every further fixed \NSi{NS-F-P162-040}{\eta}-derivative order; no
simultaneous smallness condition for all derivative orders is needed. The fixed bump shapes
also make the correction small in the radial derivatives used by the shear. Its partial radial
